% D-glossary.tex: термины в том смысле, в каком ими пользуется документ, и сокращения.
\section{Словарь}
\label{app:glossary}

Термины расположены в алфавитном порядке. За ними следуют сокращения, которыми пользуется
текст: каждое раскрыто по-русски, а для названий стандартов и организаций сказано, что они
обозначают.

\subsection*{Термины}

\begin{description}
  \item[Аттестация.] Направленное утверждение в пределах контекста о том, что одно ведомство принимает удостоверения другого, записанное строкой. Никогда не транзитивна.
  \item[Внешнее существительное.] Названная сторона за пределами репозитория, которая сделала с ним нечто, в определённую дату, с результатом, на который можно указать: единственный род записей, который принимает сводная таблица.
  \item[Двоесказательство.] Журнал, рассказывающий двум наблюдателям две разные истории. Его доказательство есть две подписанные вершины одного размера с разными корнями.
  \item[Дерево Меркла.] Криптографический отпечаток целого множества, построенный попарным хешированием вверх до одного корня; позволяет небольшим доказательством показать, что один предмет принадлежит множеству.
  \item[Доверенный поручитель.] Человек, чьё поручительство позволяет зарегистрировать того, у кого нет документов; поручительство ограничено, сверх порога соподписывается и невидимо на удостоверении.
  \item[Доказательство включения.] Короткое доказательство того, что одна запись находится в дереве Меркла с данным корнем.
  \item[Доказательство с нулевым разглашением.] Доказательство истинности утверждения, не раскрывающее ничего сверх самой этой истинности; в Полярисе принадлежность токена к эпохе.
  \item[Доказательство согласованности.] Короткое доказательство того, что меньшее дерево Меркла есть начало большего, то есть что журнал между двумя вершинами только дописывал.
  \item[Долгосрочная проверка.] Решение о том, что подпись была действительной в момент её постановки, по свидетельствам, закреплённым в этот момент, так что решение переживает позднейший вывод ключа из обращения.
  \item[Закат.] День, когда старое удостоверение перестают принимать; рассматривается как решение с нижним пределом под ним, потому что это момент, когда новое становится обязательным.
  \item[Кандидат в выпуски.] Версия \code{1.0.0-кв.7}: все условия, которые операционный договор называет для 1.0.0, выполнены, но ни один оператор, кроме автора, её не запускал.
  \item[Канонизация.] Превращение структуры в одну точную последовательность байтов (ключи упорядочены, пробельных символов нет), чтобы подпись, сделанную на одной машине, можно было проверить на другой, которая восстанавливает те же байты.
  \item[Лента отзывов.] Подписанный перечень хешей отозванных удостоверений, выданных органом, потребляемый автономно.
  \item[Множество анонимности.] Толпа, в которой прячется доказательство принадлежности: участники той эпохи, против которой оно сделано. Ничто в схеме не задаёт ему нижнего предела выше единицы, и ни один вердикт о нём не сообщает.
  \item[Мутационное учение.] Программа, которая убирает или обращает гарантию, ограничение, триггер, отказ или входы проверки и требует, чтобы тесты, якобы её доказывающие, покраснели, причём с отрицательным контролем, чтобы ноль выживших отличался от стенда, не прогнавшего ничего.
  \item[Набор якорей.] Опубликованные ключи эмитентов, которым доверяющая сторона решила доверять. Доверие в Полярисе всегда есть принадлежность к собственному набору якорей верификатора и никогда не решение, которое эмитент принимает за него.
  \item[Нуллификатор.] Значение, выведенное из секрета и области, которое позволяет верификатору обнаружить повторное использование, не связывая использований в разных областях. В Полярисе это \code{Посейдон(секрет\,||\,область\,||\,эпоха)}, открытый вход доказательства принадлежности, свой для каждой доверяющей стороны и каждой эпохи и никогда не глобальный.
  \item[Обмен слухами.] Свидетели делятся вершинами журнала, которые им показали, чтобы расщеплённый обзор обнаруживался сравнением.
  \item[Операционный договор.] Правило, принятое 13~сентября 2026 года: единица прогресса есть пережитая внешняя зависимость, изменению продукта нужна названная внешняя причина, а всё остальное уходит в лабораторию.
  \item[Осведомлённый о принуждении.] Слово, которым парадная дверь называет код принуждения, потому что механизм замечает принуждение, но не противостоит ему; более сильное слово отвергается проверкой.
  \item[Отделяемый верификатор.] Программа, проверяющая предметы Поляриса без его кода, без базы данных и без сети: пакет \code{полярис-проверка} и сценарий, ставший переходником на него.
  \item[Отказоустойчивость.] База данных, переживающая потерю ведущего узла повышением реплики по аренде, отработанная в учениях и измеренная.
  \item[Пакет подлинности.] Удостоверение в виде файла: значение токена, его подпись и открытый ключ эмитента. То, что держит держатель и проверяет посторонний.
  \item[Подписанная вершина дерева.] Подписанное обязательство журнала обо всей его истории на некотором размере.
  \item[Положительный контроль.] Та половина измерения, которая доказывает, что инструмент вообще способен что-то увидеть: противник, обязанный преуспеть на связываемом населении; изменённое дерево, на котором проверка обязана упасть; пропатченный верификатор, который набор тестов обязан заметить.
  \item[Попарный идентификатор.] Отдельный дескриптор для каждой доверяющей стороны, чтобы две стороны не могли соединить свои журналы. В Полярисе субъект входного токена и дескриптор предъявления выводятся в области доверяющей стороны; это ограничивает то, что верификатор хранит, но не то, что ему показывают.
  \item[Постквантовый.] Относящийся к криптографии, стойкость которой не опирается на задачи, эффективно решаемые квантовым вычислителем. В настоящем документе утверждение об алгоритмах и никогда о системе.
  \item[Принцип двух свидетелей.] Ни одному криптографическому вердикту не доверяют со слов одной реализации; вторая, написанная независимо, обязана согласиться или явно воздержаться.
  \item[Раздел.] Таблица, разбитая по отметке времени на помесячные части, прозрачные для запросов, чтобы старые месяцы можно было отсоединять в архив.
  \item[Резервный регион.] Второй кластер со своим хранилищем аренды, получающий поток изменений от первого асинхронно и не выбирающий собственного ведущего, пока его не повысят; его точка восстановления измерена и не нулевая.
  \item[Реплика.] Копия базы данных, следующая за ведущим узлом с некоторым отставанием; безопасна для чтения неизменяемого материала и небезопасна для решения о текущих полномочиях.
  \item[Свидетель.] Независимая сторона, которая соподписывает вершины журнала, признанные ею согласованными, и отвергает ответвление.
  \item[Сводная таблица.] Файл \code{лаборатория/ВНЕШНИЕ-СУЩЕСТВИТЕЛЬНЫЕ.мд}, который решает, продолжится ли проект, и держит лишь то, что сделали названные внешние стороны, причём ноль и пустота суть допустимые значения.
  \item[Список доверия.] Подписанное утверждение о каждом ключе, известном экземпляру, с состоянием каждого ключа и моментами, с которых эти состояния действуют.
  \item[Только на добавление.] Таблица, в которой ОБНОВИТЬ и УДАЛИТЬ отвергаются триггером базы данных, так что к истории можно лишь дописывать. Аудит как сама запись есть названный образец применения этого к каждой таблице.
  \item[Уровень достоверности.] То, что регистрация может честно утверждать о том, как человек был подтверждён; выводится из записанных рядом свидетельств и никогда не вбивается оператором.
  \item[Уровень раскрытия.] Что записывает проверка: нулевое разглашение (без идентификатора токена), избирательное (названные признаки) или полное (личность, с записью в журнал).
  \item[Утверждение о состоянии.] Короткоживущее подписанное эмитентом утверждение о текущем состоянии удостоверения, подшиваемое к предъявлению, чтобы полномочия были разрешимы автономно.
  \item[Федерация.] Независимые органы, у каждого свой корень, признающие друг друга через явные аттестации, без центральной службы.
  \item[Эпоха.] Закрытое неизменяемое обязательство о множестве действующих токенов на некоторый момент в виде корня Меркла, против которого строятся доказательства принадлежности.
\end{description}

\subsection*{Сокращения}

\begin{description}
  \item[асз.] Аудит как сама запись; сокращение в именах проверок.
  \item[В3К.] Всемирный консорциум Паутины, организация, выпускающая стандарты для веба, в том числе формат проверяемого удостоверения.
  \item[ВебАутн.] Стандарт подтверждения владения аппаратным ключом или встроенным средством устройства; операторы используют его как второй фактор после пароля.
  \item[ДЖСОН.] Текстовый формат записи структурированных данных; в нём записаны все подписываемые утверждения Поляриса.
  \item[ЕС256.] Классический алгоритм подписи на эллиптической кривой П-256, который предписывает профиль ХАИП.
  \item[ИСО/МЭК 18013-5.] Международный стандарт мобильного водительского удостоверения.
  \item[КБОР.] Двоичный формат записи структурированных данных, используемый мобильными документами по ИСО 18013-5.
  \item[кв.] Кандидат в выпуски: пометка версии, ещё не объявленной окончательной, например \code{1.0.0-кв.7}.
  \item[КСР.] Комплект средств разработки: библиотека, которую встраивает в свою программу доверяющая сторона.
  \item[куар.] Двумерный штрихкод, который считывает камера; по нему передаётся предъявление.
  \item[МЛ-ДСА.] Схема цифровой подписи на модульных решётках, стандартизованная как ФИПС 204; наборы параметров 65 и 87 принимаются. Постквантовая по своим допущениям; репозиторий заявляет гибкость, а не окончательно установленную стойкость.
  \item[МПЗ.] Межсайтовая подделка запроса: атака, при которой чужой сайт заставляет браузер отправить запрос от имени вошедшего пользователя.
  \item[НИ.] Непрерывная интеграция: автоматический прогон сборки, тестов и учений при каждой отправке изменений.
  \item[НИСТ.] Национальный институт стандартов и технологий США; его публикация СП 800-63 описывает уровни достоверности цифровой личности.
  \item[НПМ.] Публичный реестр пакетов для языка ТайпСкрипт.
  \item[НР.] Нулевое разглашение.
  \item[НФС.] Связь ближнего поля: беспроводной обмен на расстоянии нескольких сантиметров, по которому карта отвечает считывателю.
  \item[ОИД4ВП.] Протокол ОпенАйДи для проверяемых предъявлений, по которому кошелёк предъявляет удостоверение верификатору; путь, на котором с Полярисом впервые заговорил сторонний кошелёк.
  \item[ПайПИ.] Публичный реестр пакетов для языка Питон.
  \item[ПИН, ПУК.] Личный код, отпирающий карту, и код разблокировки, снимающий блокировку после исчерпания попыток ввода ПИН.
  \item[пкк.] Постквантовая криптография; сокращение в именах проверок.
  \item[ПККС\#11.] Стандартный программный интерфейс к аппаратным модулям безопасности и токенам, внутри которых вырабатываются и хранятся ключи.
  \item[ПКСЕ.] Расширение входа по коду авторизации, доказывающее владение, чтобы перехваченный код нельзя было обменять.
  \item[ППИ.] Прикладной программный интерфейс: набор запросов, которыми одна программа пользуется услугами другой.
  \item[РФС.] Серия открытых документов, в которых публикуются стандарты интернета; например, РФС 6962 описывает журналы прозрачности сертификатов.
  \item[СД-ЖВТ ПУ.] Формат удостоверения, в котором эмитент подписывает дайджесты утверждений, а держатель раскрывает лишь те, о которых попросил верификатор; формат, который предъявил первый сторонний кошелёк.
  \item[СКЛ.] Язык запросов к реляционным базам данных; на нём написаны схема, триггеры и процедуры Поляриса.
  \item[СЛХ-ДСА.] Схема цифровой подписи на одних лишь хеш-функциях, стандартизованная как ФИПС 205; запасной вариант на случай падения решёточных допущений.
  \item[СУК.] Служба управления ключами, облачная служба хранения ключей.
  \item[ТЛА+.] Язык формальных спецификаций, модели на котором перебирают все чередования событий и проверяют, что свойство держится.
  \item[ТЛС.] Протокол защищённого соединения поверх сети; в Полярисе его кромка согласует гибридный постквантовый обмен ключами.
  \item[ФИПС.] Федеральные стандарты обработки информации США; ФИПС 202 описывает хеш-функции ША-3, ФИПС 204 подпись МЛ-ДСА, ФИПС 205 подпись СЛХ-ДСА.
  \item[Х.509.] Стандартный формат сертификата открытого ключа.
  \item[ХАИП.] Профиль высокой надёжности взаимодействия для ОИД4ВП, который закрепляет формат удостоверения, алгоритм подписи и режим ответа.
  \item[ХТТП.] Протокол передачи гипертекста, по которому обмениваются запросами веб-службы.
  \item[ША3-256.] Криптографическая хеш-функция семейства ША-3 (ФИПС 202), дающая 256-битный дайджест.
\end{description}
